% Generated draft: requires mathematical and source review.
\documentclass[11pt]{article}
\usepackage[margin=25mm]{geometry}
\usepackage{fontspec}
\setmainfont{Latin Modern Roman}
\usepackage{amsmath,amssymb,mathtools,unicode-math}
\setmathfont{Latin Modern Math}
\usepackage{enumitem,needspace,xurl,hyperref}
\setlength{\parindent}{0pt}
\setlength{\parskip}{5pt}
\setlength{\emergencystretch}{3em}
\newcommand{\N}{\mathbb{N}}
\newcommand{\Z}{\mathbb{Z}}
\newcommand{\Q}{\mathbb{Q}}
\newcommand{\R}{\mathbb{R}}
\newcommand{\C}{\mathbb{C}}
\newcommand{\F}{\mathbb{F}}
\newcommand{\A}{\mathbb{A}}
\newcommand{\PP}{\mathbb{P}}
\newcommand{\E}{\mathbb{E}}
\newcommand{\eps}{\varepsilon}
\newcommand{\dd}{\,\mathrm{d}}
\begin{document}
\section*{UnsolvedMath ID: 30000308}
\textbf{Source problem code:} catalogue entry
\subsection{Definitions and mathematical statement}
Classify the algebras arising in the spectral-measure theory of free quantum groups.

The supplied assessment describes the intended objects as universal Hopf algebras. A Hopf algebra has an associative unital algebra structure, a coproduct $\Delta:A\to A\otimes A$, a counit $\epsilon:A\to\C$, and an antipode satisfying the Hopf identities; a Hopf $*$-algebra also has the relevant compatible involution. This explains terminology only: the source gives no generators, relations, universal property, quantum-group family, or base-field convention defining the class to be classified.

A precise classification would first specify a class $\mathcal C$ of the intended algebras and an equivalence relation $\sim$, and then determine $\mathcal C/\!\sim$, with a complete description of representatives and criteria for equivalence. Neither $\mathcal C$ nor $\sim$ is recoverable from the supplied statement. In particular, classification as abstract algebras, as Hopf algebras, as Hopf $*$-algebras, or as chosen $C^*$-completions are different questions. No one of them is selected here.

The phrase ``arising in the spectral-measure theory'' also lacks the defining observable, state, or association from a quantum group to a spectral measure or algebra. The supplied assessment mentions spectral examples, but includes none of them. Thus this is a broad classification request with missing object definitions, not a conjecture that a named list is complete or that spectral measures are complete invariants. No specific free orthogonal, free unitary, quantum permutation, or other family is silently substituted for the report's class.
\subsection{Short English statement}
The source asks for a classification of certain universal Hopf algebras connected with spectral measures of free quantum groups. It does not define the algebra class, the spectral construction, or the allowed notion of equivalence, so a precise classification theorem cannot yet be stated.
\subsection{Sources}
\begin{itemize}
\item Local UnsolvedMath registry, permanent ID 30000308, code catalogue entry, ``Spectral-Measure Algebras of Free Quantum Groups.'' Proposed year 2005; proposer unspecified.
\item \emph{$C^*$-Algebren} (2006), Oberwolfach Reports, pp.\ 2305--2374, DOI 10.4171/owr/2005/41; original report \url{https://doi.org/10.4171/owr/2005/41}. The displayed report year 2006 differs from the DOI and proposed year 2005.
\item Statement status exact; the supplied verification says the statement is coherent and admits no defensible textual repair. This does not provide the omitted object definitions.
\item Literature triage checked 2026-08-21: open. It calls this a fundamental broad classification problem for universal Hopf algebras and says the spectral examples in the passage do not answer it. No examples or further literature are supplied. Created 2026-08-25; updated 2026-09-16. Status and cited claims require independent verification.
\item Registry curation by Ulam AI / UnsolvedMath Contributors, CC BY 4.0; the underlying report retains its supplied provenance. Unreviewed statement-only draft.
\end{itemize}
\end{document}
